\documentclass[10pt,letterpaper]{article}

\usepackage[margin=0.65in]{geometry}
\usepackage[hidelinks]{hyperref}
\usepackage{enumitem}
\usepackage{tabularx}
\usepackage{xcolor}
\usepackage[T1]{fontenc}
\usepackage[utf8]{inputenc}
\usepackage{titlesec}
\usepackage{array}
\usepackage{parskip}

\setlength{\parindent}{0pt}
\setlength{\parskip}{0.35em}
\setlist[itemize]{leftmargin=1.2em, itemsep=0.12em, topsep=0.1em, parsep=0pt}
\setlist[description]{leftmargin=0em, itemsep=0.12em, topsep=0.1em, style=nextline}
\titlespacing*{\section}{0pt}{0.45em}{0.2em}
\titleformat{\section}{\large\bfseries}{}{0pt}{}
\pagestyle{empty}
\urlstyle{same}

\newcommand{\sectionrule}{\vspace{-0.1em}\hrule\vspace{0.25em}}
\newcommand{\entry}[4]{%
	\textbf{#1} \hfill #2\\
	\textit{#3}\nopagebreak\vspace{-0.55em}
	#4\vspace{0.15em}
}

\begin{document}
	
	{\LARGE \textbf{Andre Ricardo Goncalves}}\\[0.25em]
	Machine Learning Research Scientist\\
	\href{mailto:goncalves1@llnl.gov}{goncalves1@llnl.gov} \textbar{} +1 (925) 423-3229 \textbar{} Livermore, CA\\
	\href{https://andreric.github.io}{andreric.github.io} \textbar{} \href{https://linkedin.com/in/argoncalves}{linkedin.com/in/argoncalves} \textbar{} \href{https://scholar.google.com/citations?user=pn7hL2cAAAAJ&hl=en}{Google Scholar}
	
	\section*{Summary}
	\sectionrule
	Machine learning research scientist with 15+ years of experience developing statistical learning, optimization, and deep learning methods for biological, clinical, and scientific data. Current work spans antibody design, microbiome modeling, and large-scale clinical foundation models trained on structured electronic health record data. Published in \textit{Nature}, \textit{Science Advances}, \textit{The Lancet Digital Health}, and \textit{JMLR}.
	
	\section*{Research Experience}
	\sectionrule
	
	\entry{Lawrence Livermore National Laboratory}{Oct 2017 -- Present}{Machine Learning Research Scientist, Machine Learning Group}{
		\begin{itemize}
			\item Build and lead machine learning efforts across protein design, clinical machine learning, and microbiome research in collaboration with biologists, clinicians, and applied scientists.
			\item Develop methods for antibody developability prediction, in silico library design, AlphaFold3-based binding estimation, and simulation tools for mutagenesis and experimental design, reducing wet-lab iteration in therapeutic and biomaterials campaigns.
			\item Train and adapt billion-parameter clinical foundation models on structured event sequences at thousand-GPU scale for biosurveillance, infectious disease detection, and diagnosis prediction, and developed causal ML methods for drug repurposing in ALS.
			\item Led microbiome disease-prediction work integrating metagenomic profiles with host and protocol metadata across tens of thousands of samples.
			\item Mentored 10+ undergraduate and graduate students and postdoctoral researchers on interdisciplinary machine learning projects.
	\end{itemize}}
	
	\entry{Center for Research and Development in Telecommunication (CPqD)}{Apr 2015 -- Aug 2017}{Machine Learning Researcher}{
		\begin{itemize}
			\item Developed deep learning and machine learning models for voice biometrics and audio processing, including presentation-attack detection using spectral, temporal, and convolutional approaches.
			\item Contributed to production-oriented speaker recognition systems deployed across multiple companies in Brazil.
			\item Co-inventor on a Brazilian patent for automatic voice liveness detection.
	\end{itemize}}
	
	\entry{Doctoral and Graduate Research}{2009 -- 2016}{University of Campinas \& University of Minnesota, Twin Cities}{
		\begin{itemize}
			\item Developed sparse and structural multitask learning methods for high-dimensional prediction problems in climate science and biomedical data, as part of MS and PhD research.
			\item Conducted 1.5 years of doctoral research at the University of Minnesota (2013--2014) as a Science without Borders visiting scholar.
			\item Published on multitask learning, structured estimation, and Earth system modeling in venues including \textit{JMLR}, \textit{Computing in Science \& Engineering}, \textit{AAAI}, and \textit{CIKM}.
	\end{itemize}}
	
	\section*{Education}
	\sectionrule
	
	\textbf{Ph.D. in Electrical and Computer Engineering}, University of Campinas and University of Minnesota \hfill 2011 -- 2016\\
	Thesis: \textit{Sparse and Structural Multitask Learning}\\[0.2em]
	\textbf{M.S. in Electrical and Computer Engineering}, University of Campinas \hfill 2009 -- 2011\\
	Thesis: \textit{Online Learning in Estimation of Distribution Algorithms for Dynamic Environments}\\[0.2em]
	\textbf{B.S. in Computer Science}, State University of Londrina \hfill 2005 -- 2008\\
	Undergraduate capstone thesis: \textit{Fundamentals and Applications of Machine Learning}
	
	\section*{Honors and Awards}
	\sectionrule
	\begin{itemize}
		\item Director's Science and Technology Award recipient and Director's Award in Publication, Lawrence Livermore National Laboratory, 2022.
		\item Best Ph.D. Thesis in Computer Engineering, University of Campinas, 2016.
		\item Science without Borders scholarship for visiting research at the University of Minnesota, 2013--2014.
		\item Highest GPA in the Computer Science graduating class, State University of Londrina, 2008.
	\end{itemize}
	
	\section*{Selected Publications}
	\sectionrule
	\begin{itemize}[leftmargin=1.4em]
		\item Reimer, R. J., Soper, B., Wilson, J. L., \textbf{Goncalves, A. R.}, et al. Identification of drug repurposing candidates for amyotrophic lateral sclerosis using electronic health records: A retrospective cohort study. \textit{The Lancet Digital Health}, 2026.
		\item \textbf{Goncalves, A. R.}, Pico, J. C., Hu, Y., Schlessinger, D., Greene, J., O'Suilleabhain, L., et al. AI-enabled diagnostic prediction within electronic health records to enhance biosurveillance and early outbreak detection. \textit{Under Review}, 2026.
		\item Zhu, F., Rajan, S., Hayes, C. F., Kwong, K. Y., \textbf{Goncalves, A. R.}, et al. Preemptive optimization of a clinical antibody for broad neutralization of SARS-CoV-2 variants and robustness against viral escape. \textit{Science Advances}, 2025.
		\item Desautels, T. A., Arrildt, K. T., Zemla, A. T., Lau, E. Y., Zhu, F., Ricci, D., Cronin, S., \textbf{Goncalves, A. R.}, et al. Computationally restoring the potency of a clinical antibody against Omicron. \textit{Nature}, 2024.
		\item \textbf{Goncalves, A. R.}, Ranganathan, H., Valdes, C., Zhu, H., Zhang, B., Kok, C. R., et al. Beyond microbial abundance: Metadata integration enhances disease prediction in human microbiome studies. \textit{Frontiers in Microbiology}, 2025.
		\item \textbf{Goncalves, A. R.}, Ray, P., Soper, B., Stevens, J., Coyle, L., and Sales, A. P. Generation and evaluation of synthetic patient data. \textit{BMC Medical Research Methodology}, 2020.
		\item \textbf{Goncalves, A. R.} and Von Zuben, F. J. Multi-task sparse structure learning with Gaussian copula models. \textit{Journal of Machine Learning Research}, 2016.
	\end{itemize}
	
	Full publication list: \href{https://scholar.google.com/citations?user=pn7hL2cAAAAJ&hl=en}{Google Scholar profile}.
	
	\section*{Professional Service and Intellectual Property}
	\sectionrule
	Peer reviewer and program committee member for venues including AAAI, SDM, \textit{npj Digital Medicine}, and \textit{PLOS ONE}. Co-inventor on Brazilian patent BR 10 2015 017556-6, automatic voice liveness detection system.
	
\end{document}